\documentclass[12pt,a4paper]{article}

% Кодировка и русский язык
\usepackage[utf8]{inputenc}
\usepackage[T2A]{fontenc}
\usepackage[russian]{babel}

% Математические пакеты
\usepackage{amsmath,amssymb,amsfonts}

% Разметка страницы
\usepackage{geometry}
\geometry{
  a4paper,
  left=2cm,
  right=2cm,
  top=2.5cm,
  bottom=2.5cm
}

% Дополнительно: абзацы, межстрочный интервал при желании
\setlength{\parindent}{1.25cm}
\setlength{\parskip}{0pt}

\begin{document}

\begin{center}
    {\Large Решения задач варианта 1}\\[1em]
\end{center}

% --------------------------------------------------
% ЗДЕСЬ ВСТАВИТЬ БЛОК С РЕШЕНИЯМИ
% (скопируй из файла, который я называл answers.txt)
% --------------------------------------------------

\section*{Задача 1}

В сундуке находятся $15$ кучек красной пыли, $12$ кристаллов лазурита и $9$ изумрудов, причём на ощупь все предметы неотличимы. Нам нужно гарантировать, что среди вытянутых предметов окажется хотя бы $3$ изумруда.

Рассуждаем по принципу «наихудшего случая». Представим, что удача полностью отвернулась, и вы сначала вытягиваете все возможные не-изумруды, то есть всю красную пыль и весь лазурит. Всего таких предметов $15 + 12 = 27$, и если достать ровно $27$ предметов, может случиться так, что это окажутся только они, без единого изумруда.

Дальше, начиная с $28$-го предмета, вы уже не можете тянуть только не-изумруды, потому что они закончились: в сундуке остаются только $9$ изумрудов. Поэтому:
\begin{itemize}
  \item при вытягивании $28$ предметов у вас гарантированно окажется хотя бы $1$ изумруд;
  \item при вытягивании $29$ предметов --- хотя бы $2$ изумруда;
  \item при вытягивании $30$ предметов --- хотя бы $3$ изумруда.
\end{itemize}
Это следует из того, что после выбора $27$ не-изумрудов каждый следующий предмет обязан быть изумрудом, потому что других типов уже нет.

Таким образом, минимальное количество предметов, которые надо вытащить вслепую, чтобы среди них гарантированно было не менее трёх изумрудов, равно
\[
27 + 3 = 30.
\]

\subsection*{Ответ}
Нужно вытащить минимум $30$ предметов.

\newpage

\section*{Задача 2}

В Нижнем мире построены $13$ порталов; требуется соединить каждый портал прямой ледовой дорогой ровно с $5$ другими порталами. То есть в графе, вершинами которого являются порталы, каждая вершина должна иметь степень $5$.

Переформулируем задачу на языке графов. Имеется простой неориентированный граф с $n = 13$ вершинами, и нужно, чтобы степень каждой вершины была равна $5$. Сумма степеней всех вершин такого графа равна
\[
\sum_{v} \deg(v) = 13 \cdot 5 = 65.
\]
Но в любом неориентированном графе выполняется фундаментальное свойство: сумма степеней всех вершин равна удвоенному количеству рёбер, то есть
\[
\sum_{v} \deg(v) = 2E,
\]
где $E$ --- число рёбер. Отсюда следует, что сумма степеней всегда чётна.

В нашем случае сумма степеней рассчитана как $65$, а это нечётное число. Значит, не существует никакого неориентированного графа (а значит, и никакой схемы дорог между порталами), в котором было бы $13$ вершин и у каждой степень ровно $5$. Проблема чисто математическая и не зависит от того, можно ли пересекать дороги в Майнкрафте или строить их на разных высотах.

Следовательно, строители действительно правы: задачу выполнить невозможно.

\subsection*{Ответ}
В бланк: Да.\\
Пояснение: такой граф невозможен, так как сумма степеней $13 \cdot 5 = 65$ должна быть чётной, а она нечётна.

\newpage

\section*{Задача 3}

Имеется площадка $10 \times 10$ из обычного камня (всего $100$ блоков), каждый блок соприкасается с соседями по вертикали и горизонтали. Игрок заражает скалком ровно $9$ выбранных блоков, далее заражение происходит так: каждую минуту любой блок камня, соприкасающийся минимум с двумя заражёнными блоками, сам становится заражённым. Распространение идёт до тех пор, пока есть блоки, удовлетворяющие правилу. Нужно понять: можно ли так выбрать начальные $9$ блоков, чтобы в итоге заразился каждый из $100$ блоков.

Будем смотреть на \emph{периметр} заражённой области: это количество рёбер между заражёнными и незаражёнными (или внешней границей поля). Для каждого состояния системы периметр --- это число таких граничных рёбер.

Докажем, что при распространении заражения по данному правилу периметр никогда не увеличивается.

Рассмотрим момент, когда новая клетка $X$ становится заражённой. По условию у неё есть как минимум два заражённых соседа по сторонам. Пусть у неё $k$ заражённых соседей и $4-k$ незаражённых (или внешняя граница). До заражения клетки $X$ каждый её сосед-заражённый даёт одно граничное ребро: между заражённой клеткой и незаражённой $X$. Таких рёбер ровно $k$, они входят в периметр.

После того как клетка $X$ заражается, эти $k$ рёбер перестают быть граничными: теперь по этим сторонам заражённая клетка соприкасается с заражёнными соседями, и периметр по ним уменьшается на $k$. Зато появляются новые граничные рёбра по сторонам, где у $X$ незаражённые соседи или внешняя граница: их ровно $4-k$. Таким образом, изменение периметра при заражении одной клетки равно
\[
\Delta P = (4-k) - k = 4 - 2k.
\]

Но по условию $k \ge 2$ (у новой заражаемой клетки минимум два заражённых соседа). Значит
\[
\Delta P = 4 - 2k \le 4 - 4 = 0.
\]
То есть при каждом акте заражения периметр либо уменьшается, либо остаётся прежним, но \emph{никогда не увеличивается}.

Теперь оценим начальный периметр и периметр всей площадки. Если игрок выбирает начальные $9$ заражённых клеток так, чтобы они образовывали максимально компактную область (например, квадрат $3 \times 3$ внутри поля), то периметр этой области будет равен
\[
P_{\text{нач}} = 4 \cdot 3 = 12 \quad \text{по одной стороне} \quad \Rightarrow \quad P_{\text{нач}} = 12 + 12 = 24,
\]
или, более точно, для стандартной конфигурации можно добиться периметра $36$ (каждая из $9$ клеток даёт в среднем по $4$ сторон, но внутренние стороны не учитываются в периметре; в любом случае важно, что начальный периметр меньше периметра всей площадки). Для всей площадки $10 \times 10$ периметр равен
\[
P_{\text{вся}} = 4 \cdot 10 = 40.
\]

Так как при каждом шаге заражения периметр никогда не возрастает, начальный периметр заражённой области (который можно сделать равным $36$) не может увеличиться до $40$, который у всей площадки. Следовательно, конфигурация, в которой заражены все $100$ клеток, недостижима.

Таким образом, при любом выборе начальных $9$ заражённых блоков скалк не сможет поглотить всю площадку $10\times 10$.

\subsection*{Ответ}
В бланк: Нет.\\
Пояснение: при данном правиле заражения периметр заражённой области не может увеличиваться, а периметр стартовой области из $9$ клеток меньше периметра всей площадки ($36 < 40$), поэтому заразить все $100$ клеток невозможно.
\newpage

\section*{Задача 4}

Есть $5000$ пронумерованных ламп, каждая изначально выключена. В тик $k$ переключаются все лампы с номерами, кратными $k$ (то есть состояние каждой такой лампы меняется: «0» $\leftrightarrow$ «1»). Нужно ответить на два вопроса:
\begin{enumerate}
  \item Сколько ламп останется во включённом состоянии?
  \item Какой номер у самой последней включённой лампы?
\end{enumerate}

Рассмотрим произвольную лампу с номером $n$ (от $1$ до $5000$). Она переключается в те тики $k$, для которых $n$ кратно $k$, то есть $k$ является делителем числа $n$. Значит, количество переключений у лампы с номером $n$ равно количеству делителей числа $n$. Изначально лампа выключена, поэтому она окажется включённой тогда и только тогда, когда число её переключений нечётно.

Из теории чисел известно, что у числа нечётное количество делителей только в одном случае: если это число является квадратом целого числа. Причина проста: делители обычно разбиваются на пары $d$ и $n/d$, а у квадратного числа $n = m^2$ один из делителей (именно $m$) является «сам себе парой», поэтому общее количество делителей становится нечётным. Следовательно, включёнными останутся лампы, номера которых являются полными квадратами.

Теперь нужно посчитать, сколько таких квадратов лежит между $1$ и $5000$ включительно. Пусть $k^2 \leq 5000$. Максимальное целое $k$, удовлетворяющее этому, равно
\[
k = \left\lfloor \sqrt{5000} \right\rfloor.
\]
Заметим, что $70^2 = 4900$, а $71^2 = 5041 > 5000$, значит $k = 70$. Таким образом, подходящие квадраты --- это $1^2, 2^2, \dots, 70^2$, всего $70$ ламп.

Самая последняя включённая лампа --- это лампа с наибольшим номером среди квадратов, не превосходящих $5000$, то есть с номером
\[
70^2 = 4900.
\]

\subsection*{Ответ}
В бланк: $70\ 4900$.\\
Пояснение: включёнными останутся лампы с номерами $1^2, 2^2, \dots, 70^2$; их ровно $70$, наибольший номер --- $4900$.

\newpage

\section*{Задача 5}

Арена --- прямоугольник $1920$ на $1080$, углы имеют координаты $(0,0)$, $(1920,0)$, $(0,1080)$, $(1920,1080)$. Из точки $(0,0)$ выпускается стрела под углом $45^\circ$, то есть по диагонали, и при отражении от стен угол падения равен углу отражения. Стрела летит, пока не попадёт точно в один из четырёх углов-мишеней.

Нас просят определить:
\begin{itemize}
  \item в какой из трёх остальных углов она попадёт;
  \item сколько раз она отскочит от стен до попадания (само попадание в мишень не считается отскоком).
\end{itemize}

Стандартный приём для таких задач с отражениями --- «разворачивать» прямоугольник, вместо того чтобы отслеживать отражения. То есть мы мысленно продолжаем движение стрелы по прямой линии, а каждый раз, когда она должна была отразиться, рассматриваем её как будто она вошла в соседний прямоугольник, являющийся зеркальным отражением исходного. В этой развёрнутой картине стрела просто летит по прямой под углом $45^\circ$ от исходной точки $(0,0)$ в бесконечной решётке прямоугольников.

Движение под углом $45^\circ$ означает, что при некотором масштабе координаты вдоль осей растут одинаково: траектория имеет вид $y = x$ (если считать масштаб одинаковым по осям). Углы исходного прямоугольника в развёрнутой решётке будут располагаться в точках вида $(k \cdot 1920,\ \ell \cdot 1080)$, где $k$ и $\ell$ --- целые числа. Попасть в угол означает, что траектория $y = x$ совпадёт с одной из таких точек:
\[
x = k \cdot 1920, \quad y = \ell \cdot 1080, \quad \text{и при этом } y = x.
\]
То есть нужно, чтобы выполнялось
\[
k \cdot 1920 = \ell \cdot 1080.
\]

Сократим это равенство. Найдём $\gcd(1920, 1080)$. Разложим числа:
\[
1920 = 30 \cdot 64,\qquad 1080 = 30 \cdot 36.
\]
Равенство $k\cdot 1920 = \ell\cdot 1080$ переписывается как
\[
k \cdot 64 = \ell \cdot 36.
\]
У чисел $64$ и $36$ наибольший общий делитель равен $4$, поэтому можно представить
\[
k = 9t,\quad \ell = 16t,
\]
где $t$ --- целое число. Нас интересует первый момент попадания в угол, то есть минимальное положительное $t$ --- это $t=1$. Тогда
\[
k = 9,\quad \ell = 16,
\]
и точка попадания в развёрнутой решётке:
\[
x = k \cdot 1920 = 9 \cdot 1920 = 17280, \quad y = \ell \cdot 1080 = 16 \cdot 1080 = 17280.
\]

За время движения до этой точки стрела пересекает:
\begin{itemize}
  \item вертикальные границы --- $9$ раз (каждая соответствует отражению от вертикальной стены),
  \item горизонтальные границы --- $16$ раз (каждая соответствует отражению от горизонтальной стены).
\end{itemize}
Всего она $9 + 16 = 25$ раз пересечёт границы, то есть $25$ раз отскочит от стен.

Какой именно угол исходного прямоугольника соответствует этой точке, определяется чётностью числа «отзеркаливаний» по каждой оси. В нашем случае после соответствующих отражений траектория приводит стрелу в правый верхний угол $(1920,1080)$ исходной арены.

\subsection*{Ответ}
В бланк: $(1920,1080)\ 25$.\\
Пояснение: стрела, двигаясь по диагонали и отражаясь от стен, эквивалентно летит по прямой $y = x$ в развёрнутой решётке прямоугольников, впервые попадает в вершину, соответствующую правому верхнему углу, после $25$ пересечений границ.

\end{document}